\chapter{Foundations: From Learning Problems to Multilayer Perceptrons}

Machine intelligence starts from a simple observation: the system designer usually does not know the correct input--output rule in advance. The task is therefore to specify a model family, define a criterion of success, and learn parameters from data rather than write the full rule by hand.

\section{Learning as model selection}

At a high level, a learning problem has three ingredients:
\begin{enumerate}
  \item a hypothesis class,
  \item a loss function,
  \item and an optimization procedure.
\end{enumerate}

\begin{keyequation}[Learning pipeline]
\[
\text{data} \;\longrightarrow\; \text{model class} \;\longrightarrow\; \text{loss} \;\longrightarrow\; \text{optimization}
\]
\tcblower
Learning is not ``fit a curve'' in the abstract. It is the joint choice of representation, objective, and search procedure.
\end{keyequation}

This viewpoint already separates the main paradigms of the course:
\begin{itemize}
  \item \textbf{supervised learning}: learn from labeled pairs $(\vec{x}^{(\alpha)}, y^{(\alpha)})$,
  \item \textbf{unsupervised learning}: discover structure in $\vec{x}^{(\alpha)}$ alone,
  \item \textbf{reinforcement learning}: learn from trajectories with delayed reward.
\end{itemize}

\begin{aside}[Core intuition]
The paradigms differ mainly by what feedback is available. Labels give direct error signals, unsupervised objectives invent structure from the inputs themselves, and reinforcement learning must assign credit across time.
\end{aside}

\section{Connectionist neurons}

The elementary computational unit of early machine learning is the connectionist neuron,
\[
  y_i = f\!\left(\sum_j w_{ij} x_j - \theta_i\right).
\]
Here $\vec{x}$ is the input, $\wvec_i$ is a learned weight vector, $\theta_i$ is a threshold or bias term, and $f$ is a transfer function such as a threshold, sigmoid, $\tanh$, or ReLU.

Two interpretations matter:
\begin{enumerate}
  \item \textbf{feature detector}: $\wvec_i$ selects a preferred direction in input space,
  \item \textbf{decision unit}: the sign of $\wvec_i^\top \vec{x} - \theta_i$ tells us on which side of a hyperplane the input lies.
\end{enumerate}

\begin{figure}[t]
\centering
\begin{tikzpicture}[scale=0.95, >=Latex]
  \draw[->] (0,0) -- (5.8,0) node[right] {$x_1$};
  \draw[->] (0,0) -- (0,4.5) node[above] {$x_2$};
  \draw[very thick] (0.8,4.0) -- (4.8,0.35);
  \draw[->, ultra thick] (2.3,1.9) -- (4.1,3.0) node[right] {$\wvec$};
  \node at (4.2,3.6) {decision boundary};
  \fill (1.0,0.9) circle (2pt);
  \fill (1.6,1.3) circle (2pt);
  \fill (2.2,0.7) circle (2pt);
  \draw[fill=white] (4.0,2.8) circle (2.2pt);
  \draw[fill=white] (4.5,2.2) circle (2.2pt);
  \draw[fill=white] (3.5,3.1) circle (2.2pt);
\end{tikzpicture}
\caption{A single neuron implements a linear decision boundary. The weight vector is normal to the separating hyperplane.}
\end{figure}

The geometric takeaway is immediate: changing $\theta_i$ shifts the hyperplane, while changing $\wvec_i$ rotates it. A single neuron can therefore separate only linearly separable classes.

\section{Why multilayer perceptrons are needed}

One layer is not enough for problems such as XOR. A multilayer perceptron (MLP) composes affine maps and nonlinearities,
\[
  \hvec^{(\ell)} = \sigma\!\bigl(\Wmat^{(\ell)} \hvec^{(\ell-1)} + \vec{b}^{(\ell)}\bigr),
\]
so each layer transforms the representation before the next layer acts on it.

\begin{aside}[Representation view]
Depth is useful because later layers are not solving the original problem directly. They are solving a problem in a learned feature space created by earlier layers.
\end{aside}

Feedforward networks propagate information from input to output. Recurrent networks reuse internal state across time and are treated later in the course.

\section{Backpropagation}

The practical difficulty with MLPs is not representation but optimization. If a network produces loss $L$, training needs the derivative of $L$ with respect to every weight. Backpropagation solves this with repeated use of the chain rule.

\begin{keyequation}[Backpropagation recursion]
\[
\delta_j^{(\ell)}
=
\sigma'\!\bigl(a_j^{(\ell)}\bigr)
\sum_k w_{kj}^{(\ell+1)} \delta_k^{(\ell+1)},
\qquad
\pd{L}{w_{ji}^{(\ell)}} = \delta_j^{(\ell)} h_i^{(\ell-1)}.
\]
\tcblower
The error signal at one layer is a weighted sum of downstream error signals, rescaled by the local slope of the nonlinearity.
\end{keyequation}

The standard training loop is therefore:
\begin{enumerate}
  \item compute activations in a forward pass,
  \item evaluate the loss,
  \item propagate error signals backward,
  \item update weights using a gradient step.
\end{enumerate}

The most important intuition is local: every weight update asks, ``if this parameter changed slightly, how would the final loss change?'' Backpropagation makes that question tractable in large layered systems.
